\section{Causal reasoning checklists need matched errors and honest denominators}
\label{exec21:check-scope}
A bank may ask an agent to turn a servicing problem into a causal diagram, identify missing assumptions and explain whether an intervention is justified. A correct final answer on one familiar question is weak evidence for that workflow: the model might recognize the story while missing the mechanism. CausalityCheck, by Li, Duo, Gao and Su, offers a useful response by testing several tasks and perturbations around a seed question. The inspected original is the 27-page July 2026 ACL Findings article. Its contribution is a textual evaluation framework; it does not establish an identified banking effect or a qualified diagram-construction service.~\citep{exec21_causalitycheck}

\subsection{What the checklist tests, and what generates its gold answers}
\label{exec21:check-design}
The framework crosses four tasks---original questions, answerability judgment, process judgment and causal-validity assessment---with four settings: original problem solving, rewritten understanding, irrelevant disturbances and virtual-vocabulary reasoning. Answerability tasks remove key information; process tasks alter a proposed chain; causal-validity tasks distinguish four categories combining correctness and reliability. Virtual reasoning replaces real entities with fictional names. This design is attractive for testing whether a system attends to the provided context and whether its answer survives intended meaning-preserving changes.~\citep{exec21_causalitycheck}

The source of the labels matters. The paper uses GPT-5 to expand CausalProbe seeds and reports two rounds of manual checks with ten trained graduate students. One seed produces 36 instances across the 16 task-by-setting groups because the tasks have different numbers of answer variants. The reported sequence is 1,200 seeds, 43,200 generated instances, a selected 25\% subset of 10,800, and 8,920 retained instances after filtering. The printed cell counts reproduce that final total. These are the authors' reported construction and project-checked arithmetic, not a reproduction of the original model calls or manual judgments.~\citep{exec21_causalitycheck,exec21_causalitycheck_reference}

The four task totals in the printed figure are 482, 2,052, 2,190 and 4,196. Consequently, equal weighting of tasks or cells differs from weighting all retained instances equally. With cell accuracies $r_j$ and cell sizes $n_j$, the two summaries are
\begin{equation}
 R_{\rm macro}=\frac{1}{J}\sum_{j=1}^J r_j,
 \qquad
 R_{\rm micro}=\frac{\sum_{j=1}^J n_j r_j}{\sum_{j=1}^J n_j}.
 \label{exec21:check-weights}
\end{equation}
The paper describes its table entries as averages of related units. A product evaluation must specify the unit and aggregation weights, retain seed identifiers, and avoid treating multiple transformations of the same seed as independent observations. A reported model ranking is otherwise inseparable from the chosen mixture of tasks. Recovering raw instances and their grouping is a prerequisite for a new statistical comparison.~\citep{exec21_causalitycheck,exec21_causalitycheck_reference}

The paper's annotation table also supports a more qualified conclusion than blanket excellent agreement. It prints answerability kappa values of 0.72 and 0.78 in two settings and causal-validity values from 0.59 to 0.65. Several other entries exceed 0.8. Agreement is informative about labeling consistency, but does not make a causal interpretation true. For a bank, gold labels should distinguish claims entailed by a stipulated mechanism from domain judgments needing an owner or additional evidence. Disputed reliability categories should remain visible rather than be converted silently into certain causal facts.~\citep{exec21_causalitycheck}

\subsection{A difference of accuracies is not a frequency of reasoning errors}
\label{exec21:check-paired}
The article uses the difference between answerability and process-judgment accuracies to diagnose causal-chain errors. Table 1 reports 82.27\% and 59.46\% respectively for one model, producing a 22.81 percentage-point difference. A useful diagnostic contrast is not automatically the probability that a model gets answerability right but the chain wrong. To see the distinction, consider an ideal matched population of problems and let $A$ denote a correct answerability judgment and $B$ a correct process judgment. Then
\begin{equation}
 P(A)-P(B)=P(A\cap B^c)-P(A^c\cap B).
 \label{exec21:check-difference}
\end{equation}
The subtraction nets two different kinds of discordance. Without the joint outcomes, it does not identify either one.~\citep{exec21_causalitycheck,exec21_causalitycheck_reference}

For marginal probabilities $a=P(A)$ and $b=P(B)$, the compatible matched-error range is
\begin{equation}
 \max(0,a-b)\leq P(A\cap B^c)\leq\min(a,1-b).
 \label{exec21:check-bounds}
\end{equation}
Using $a=0.8227$ and $b=0.5946$ as an authored matched-design illustration gives a range from 22.81\% to 40.54\%. At the lower endpoint, the joint probabilities for both-correct, $A$-only, $B$-only and neither are $0.5946,0.2281,0,0.1773$. At the upper endpoint they are $0.4173,0.4054,0.1773,0$. Both reproduce the same marginals. The exact reference checks nonnegativity and every marginal. Because the paper's task populations and counts differ, these illustrative bounds are not asserted to be empirical bounds for its actual model responses. Matching must be established first.~\citep{exec21_causalitycheck_reference}

A negative accuracy difference is equally important. With matched marginals 71.68\% and 76.80\%, their difference is $-5.12$ points, but $P(A\cap B^c)$ can still range from zero to 23.2\%. The article's negative challenge score therefore does not by itself demonstrate an absence of chain errors. A proper evaluation should retain paired predictions and gold judgments at seed level and report both discordance directions, uncertainty under the seed sampling law and the intended causal level of each task. Even a well-measured process-label error rate is not a rate of incorrect interventional conclusions unless the labels establish that connection.~\citep{exec21_causalitycheck,exec21_causalitycheck_reference}

\subsection{Filtering changes the decision population}
\label{exec21:check-selection}
The paper removes incomplete fields, ambiguous relationships and sensitive material, and reports 1,880 exclusions from the 10,800 selected instances. Such filtering may be appropriate for a clearly defined benchmark, but the resulting accuracy concerns accepted instances. A production agent will also encounter underspecified, disputed and unanswerable requests. Its refusal policy, escalation burden and error rate on those requests matter to the bank. A benchmark should therefore publish the sampling and exclusion flow, reasons by category and the scope of the reported denominator. The original article's filter is not evidence that the excluded requests are solved.~\citep{exec21_causalitycheck,exec21_causalitycheck_reference}

An authored illustration makes the economic issue explicit. Start with 100 requests, retain 50 and answer 45 of those correctly. Selected accuracy is 90\%. If nothing is known about the excluded half, full-population accuracy lies between 45\% and 95\%. Suppose a correct action yields USD 10, an incorrect action loses USD 30, and every attempted action costs USD 1. At success probability $p$, stipulated expected margin is
\begin{equation}
 V(p)=10p-30(1-p)-1=40p-31.
 \label{exec21:check-value}
\end{equation}
Applying selected accuracy to every request gives USD 5, whereas the compatible full-population margins range from USD -13 to USD 7. These are sensitivity bounds under a deliberately incomplete observation law, not a confidence interval or a measured bank return. If the system abstains on excluded cases, the evaluation must instead model that action's cost and residual human workload explicitly.~\citep{exec21_causalitycheck_reference}

The same lesson applies to source filtering and data-generation quality. GPT-generated items followed by manual review may be useful, but a competing model outperforming the generator does not prove that generator-specific artifacts are absent. The paper's Appendix G offers such a relative-performance argument. A stronger design would compare independently generated and human-authored items, blind the adjudicators to model identity, retain difficult disagreements, and separate development from evaluation seeds. A generator can bias every model while still failing to rank first.~\citep{exec21_causalitycheck}

\subsection{Entity replacement is a stress test, not a unique causal diagnosis}
\label{exec21:check-rename}
The article reports higher performance on virtual-vocabulary tasks and discusses a case where replacing a familiar entity changes a model's handling of a negative conclusion. This motivates a practical perturbation test. It does not uniquely identify memorization or any single internal failure mechanism: replacement can also change tokenization, semantic cues, question difficulty and response-policy triggers. Repeated identical answers to one prompt support repeatability under that setting, not the exclusion of those alternative explanations.~\citep{exec21_causalitycheck}

For a banking prototype, the controlled object should be a known mechanism with explicit variables, action clock and admissible observations. Multiple names and paraphrases can be assigned to the same mechanism while keeping its causal query fixed. A second axis can change the mechanism while retaining familiar surface vocabulary. Contrasting these axes tests whether the system follows the stipulated structure rather than a remembered story. Labels should include answerable, unidentified, unsupported and internally inconsistent cases. This is a proposed experiment motivated by the paper, not an experiment that the library has already run.

\subsection{How to use the framework in the proposal}
\label{exec21:check-integration}
CausalityCheck's task matrix is a candidate design component for evaluating business-question translation, assumption elicitation and diagram explanation. Its textual gold labels should not become an oracle for graph identification, counterfactual effects or hidden-confounder exclusion. The article itself acknowledges limited treatment of structural rigor and counterfactual reasoning. The proposal should combine checklist perturbations with exact finite causal models, known nonidentification pairs and independently adjudicated domain tasks. That combination tests both linguistic robustness and whether the requested effect follows from the available information.~\citep{exec21_causalitycheck}

Before adopting the package or dataset, recover a pinned repository and data revision, trace the transformation and filtering code, inspect license and redistribution terms, and bind item-level provenance. The current review has inspected the genuine article, not executed its code or live models. The library, CLI and MCP should expose seed and variant identity, requested causal level, gold-label origin, inclusion status, paired predictions and uncertainty scope. Agents may organize and explain that evidence; a single aggregate accuracy must not authorize a banking intervention. Public mechanism tests can establish whether a prototype merits an expensive institutional trial while preserving the limits of textual benchmark evidence.
